\section{A known exceptional realization and further center restrictions}
\label{sec:m3-creative-20261022}

This continuation combines a realization from a known exceptional
$SU(3)$ module category with necessary restrictions on the remaining
centers and Morita constructions. A finite witness for selected center
equations is distinguished from a categorical construction throughout.
The exact source identities, proofs and replay programs are collected in
\path{evidence/rank8_creative_20261022/}.

\subsection{R982 is the neutral component of the known $E_{21}$ category}
\label{subsec:creative-e21-model}

\begin{theorem}[Exceptional $SU(3)$ realization]
\label{thm:creative-e21-model}
The ring $\sha{R8_3_982}$ is categorified by the triality-zero component
of the known unitary $E_{21}=\mathcal E^{(24)}$ chiral fusion category
associated to the conformal embedding $SU(3)_{21}\subset(E_7)_1$.
It therefore left the unresolved list at this historical checkpoint.
The exact-multiplicity-three counts at that checkpoint were $3506$ general exclusions, $16$ unitary models and
$5$ unresolved cases; cumulatively through multiplicity three they were
$4402$, $93$ and $5$, respectively. The later exclusion of $\sha{R8_3_3472}$ in Theorem~\ref{thm:r3472-general-exclusion} gives the current totals in Theorem~\ref{thm:m3-current}.
\end{theorem}
\begin{proof}
Let $\mathcal C=\operatorname{Rep}(SU(3)_{21})$ and let $A$ be the
connected commutative algebra supplied by the local conformal embedding.
Its vacuum branching object is
\[
\begin{split}
 A={}&(0,0)+(4,4)+(6,6)+(10,10)+(21,0)+(0,21)\\
    &+(13,4)+(4,13)+(10,1)+(1,10)+(9,6)+(6,9).
\end{split}
\]
The category $\mathcal E=\mathcal C_A$ of all right modules, with
relative tensor product, is a unitary fusion category. This categorical
existence is classical: B\"ockenhauer--Evans,
\emph{Modular Invariants, Graphs and $\alpha$-Induction for Nets of
Subfactors III}, Section~6.4
(\href{https://arxiv.org/abs/hep-th/9812110}{hep-th/9812110}),
construct the twenty-four irreducible chiral sectors inside a unitary
endomorphism category and give their fundamental fusion graph in
Figure~13. This is the $E_{21}$ graph, not merely an abstract
nonnegative matrix representation. The algebra/module description is
also covered by Davydov--M\"uger--Nikshych--Ostrik,
\emph{The Witt group of non-degenerate braided fusion categories},
Sections~3.3 and~6.1
(\href{https://arxiv.org/abs/1009.2117}{arXiv:1009.2117}).
The local-module category has rank two; here we use all modules.

Every summand of $A$ has triality $a-b=0\pmod3$. Decomposing the
underlying object of an $A$-module by triality decomposes its module
structure, so each simple module is homogeneous. Relative tensor
products add degrees and duals negate them. Free induction is nonzero
on each of the three source degrees, so this grading is faithful.
Its degree-zero part $\mathcal D=\mathcal E_0$ is therefore a full
unitary fusion subcategory.

For exact identification, use the twenty-four vertices and directed
fundamental adjacency matrix $G$ of Figure~43 in Coquereaux--Isasi--Schieber,
\emph{Notes on TQFT wire models and coherence equations for $SU(3)$
triangular cells}
(\href{https://arxiv.org/abs/1007.0721}{arXiv:1007.0721}).
The matrices $G,G^{\mathsf T}$ commute. Evaluation of polynomials in
these matrices at the specified unit has successive ranks
$1,3,6,10,15,21,24$ through total degree six. Thus the unit is cyclic,
and the graph determines the full fusion algebra uniquely: a polynomial
representing any basis vector at the unit is its multiplication matrix.
The exact reconstruction verifies all twenty-four integral nonnegative
fusion matrices and their associativity. Categorical existence is supplied
by the preceding construction, independently of these ring tests.

The three graph colors agree with the categorical triality, since the
unit has color zero and the fundamental graph advances color by one.
The neutral component has eight vertices. In the vertex labels of
Figure~43, the unit-preserving identification from the canonical source
basis is
\[
 (X_0,X_1,\ldots,X_7)\longmapsto
 (1_1,6_1,6_2,7_1,7_2,3_1,1_2,3_2).
\]
Every one of the $512$ neutral structure constants agrees under this
map with the exact canonical tensor of $\sha{R8_3_982}$.
The neutral multiplicity is three. This proves the required based-ring
identity, not just agreement of dimensions or spectra.
\end{proof}

As an independent scalar check, the faithful grading and the local-module
category $\operatorname{Rep}((E_7)_1)$ of dimension two give
\[
 \FPdim(\mathcal D)
 =\frac{\sqrt{2\FPdim(\mathcal C)}}3
 =\frac{\sqrt6}{\sin^2(\pi/24)\sin(\pi/12)}
 =144+96\sqrt2+80\sqrt3+56\sqrt6.
\]
The source graph and every canonical tensor entry are bound in
\path{r982/e21_model.json}; the categorical audit is
\path{unified/E21_MODEL_AUDIT.md}. Neither braiding of $\mathcal E$
nor braiding of its neutral component is asserted. The chiral systems
have a relative braiding, while their local intersection is a different
subcategory.

The construction also realizes a previously necessary algebra object.
For the simple fundamental object $V\in\mathcal E_1$, the exact tensor
identity is $V\otimes V^*=1+X_5$. Its canonical endomorphism algebra
is connected and separable, and its multiplication stays in
$\mathcal D$. Thus the earlier surviving object $1+X_5$ does carry an
algebra structure in this model. The degree-one and degree-two components give invertible
$\mathcal D$-bimodule categories with dual equivalent to $\mathcal D$.
Their exact action tensors match the two nonregular based-module types
from the earlier complete classification; thus all three necessary
based-module types are realized. Their nonregularity distinguishes both
from the identity bimodule, and the faithful $C_3$ grading supplies an
order-three subgroup of $\operatorname{BrPic}(\mathcal D)$. The graded-component module identifications and tensor-product equivalences follow from Etingof--Nikshych--Ostrik, \emph{Fusion categories and homotopy theory}, Theorem~6.1 (\href{https://arxiv.org/abs/0909.3140}{arXiv:0909.3140}). Neither
uniqueness of module associators nor the whole Brauer--Picard group is
claimed.


\subsection{Saturated dimension characters control every twist fibre}
\label{subsec:creative-saturated-conductor}

Let $\rho$ be the canonical normalized linear modular representation of
the Drinfeld center of a spherical fusion category. Its congruence
property makes it a representation of $\mathrm{SL}_2(\Z/N\Z)$ for a
suitable $N$. Write $s,t$ for the standard modular generators, and let
$P$ project onto the $t$-fixed space.

\begin{lemma}[Equality in the local conductor bound]
\label{lem:creative-local-conductor}
Let $\pi$ be a unitary representation of
$\mathrm{SL}_2(\Z/p^a\Z)$, and let $v$ be fixed by the upper
unipotent subgroup and transform by the character $\alpha$ under the
diagonal torus. If
\[
 \|P\pi(s)v\|=\operatorname{cond}(\alpha)^{-1/2}\|v\|,
\]
then the subrepresentation generated by $v$ factors through
$\mathrm{SL}_2(\Z/\operatorname{cond}(\alpha)\Z)$.
For trivial $\alpha$, this means that the generated representation is
trivial.
\end{lemma}
\begin{proof}
Apply the local norm bound of
\cite[Lemma~4.1]{SchopierayDimensionBlock} to each irreducible
summand supporting $v$. Equality holds on every such summand.
An irreducible with a nonzero unipotent-fixed vector embeds into the
permutation representation on primitive columns $(x,y)$ modulo $p^a$.
For a nontrivial character of conductor $p^c$, the full maximal
singular space of the compression $P\pi(s)P$ is spanned by
\[
 F_\alpha(x,y)=\alpha(y)[y\text{ is a unit}],\qquad
 H_\alpha(x,y)=\overline{\alpha}(x)[y\equiv0\pmod{p^c}].
\]
Indeed the large-cell line and the small-cell space are interchanged
by the compression; the small-to-large map has rank one, norm
$p^{-c/2}$, and adjoint image the displayed $H_\alpha$ line.
The two image spaces are orthogonal. Thus there are no additional
maximal singular vectors. Both displayed functions inflate from
primitive columns modulo $p^c$. That inflation subspace is invariant,
so irreducibility forces every supporting summand into it.

For trivial $\alpha$, equality of norms says that $\pi(s)v$ is
unipotent-fixed. Thus $v$ is fixed both by the upper and lower
elementary unipotent subgroups, which generate the local special
linear group.
\end{proof}

The Chinese remainder decomposition and the tensor-product
decomposition of irreducible representations imply the corresponding
global result: equality for a torus character $\chi$ forces the cyclic
representation generated by its vector to have level dividing
$f=\operatorname{cond}(\chi)$. This removes both excess prime powers
and primes not dividing $f$; no twist-order cutoff is imposed.

\begin{theorem}[Saturated conductor support]
\label{thm:creative-conductor-support}
Use the dimension-block notation
$K=\Q(\dim X:X\text{ simple})$, $D=\dim\mathcal C$, and
\[
 w_\chi=\sum_{\sigma\in\operatorname{Gal}(K/\Q)}
          \chi(\sigma)e_{A_\sigma},\qquad
 h_\chi=\sum_\sigma\frac{\chi(\sigma)}{\sigma(D)}.
\]
If the individual bound $f|h_\chi|^2=1$ is saturated, every central
simple $Z$ satisfies
\[
 \theta_Z^f\ne1\quad\Longrightarrow\quad
 c_\chi(Z):=\sum_\sigma\chi(\sigma)
                \sigma\!\left(\frac{\dim Z}{D}\right)=0.
\]
\end{theorem}
\begin{proof}
The dimension-block identities
\cite[Lemma~3.1]{SchopierayDimensionBlock} give
$\|w_\chi\|^2=|\operatorname{Gal}(K/\Q)|$ and the matrix
coefficient $|\operatorname{Gal}(K/\Q)|h_\chi$.
Saturation forces equality both in Cauchy--Schwarz and in the sharp
compression bound. Lemma~\ref{lem:creative-local-conductor} and its
Chinese remainder consequence therefore give
$\rho(t)^f\rho(s)w_\chi=\rho(s)w_\chi$.
The $Z$ coordinate of $\rho(s)w_\chi$ is $c_\chi(Z)$, whereas
$\rho(t)$ is diagonal with entry $\theta_Z$.
\end{proof}

There are useful signed refinements. For a real quadratic character
modulo a prime $p\equiv1\pmod4$, write
$h_\chi=\epsilon/\sqrt p$ with $\epsilon\in\{1,-1\}$.
Only twist $1$ and primitive twists $\exp(2\pi ik/p)$ with
$\left(\frac{k}{p}\right)=\epsilon$ can have $c_\chi(Z)\ne0$.
The exact squared masses in these fibres are
\begin{equation}\label{eq:creative-prime-masses}
 \sum_{\theta_Z=1}|c_\chi(Z)|^2=\frac{\|w_\chi\|^2}{p},
 \qquad
 \sum_{\theta_Z=e^{2\pi ik/p}}|c_\chi(Z)|^2
       =\frac{2\|w_\chi\|^2}{p}.
\end{equation}
To prove both assertions, in the primitive-column model use the
compression eigenvector $F+\epsilon\sqrt pH$.
On each large-cell unipotent orbit its $s$-image is
$\chi(y)(\chi(t)+\epsilon\sqrt p[t=0])$.
The nonzero-frequency Fourier coefficients are
$\sqrt p(\chi(k)+\epsilon)$. The total squared norm is
$2p(p-1)$, the zero-frequency squared norm is $2(p-1)$, and
each allowed nonzero-frequency squared norm is $4(p-1)$.
Equivariant isometric embeddings and orthogonal decomposition transfer
these ratios to the center representation.

For the quartic field $\Q(\sqrt2,\sqrt3)$ shared by R928, R941 and
R982, put
\[
 D=144+96\sqrt2+80\sqrt3+56\sqrt6,
 \quad \dim Z=a+b\sqrt2+c\sqrt3+d\sqrt6,
\]
and define
\[
 L_{12}=3c-a-b,\qquad L_{24}=6d-a-2b.
\]
The two saturated coordinates are $L_{12}/\sqrt{12}$ and
$L_{24}/\sqrt{24}$. Their full signed spectral supports are obtained
by tensoring the local equality vectors at conductors $4,3$ and
$8,3$. At conductor $12$, the allowed exponents are
\[
 \{0,2,3,4,5,6,7,8,9,10\}\pmod{12};
\]
at conductor $24$, they are all exponents except
$1,5,19,23\pmod{24}$. These restrictions apply only when the
corresponding mask is nonzero. The common fibre capacities are
\begin{equation}\label{eq:creative-quartic-masses}
 \sum_{\theta_Z=1}L_{12}^2=\sum_{\theta_Z=-1}L_{12}^2=4,
 \qquad
 \sum_{\theta_Z=\eta}L_{24}^2=4
       \quad(\eta\in\{1,-1,i,-i\}),
\end{equation}
and each simple satisfies $|L_{12}|,|L_{24}|\leq4$.
The local spectral distributions and their zero mixed projector
pairings are checked in the exact cyclotomic fields at orders $3,4,8$.
The two global endpoint spectra then determine every possible
distribution by a convex combination. This proves the capacities and
bounds without assuming that the two endpoint lines generate distinct
irreducible representations.

Combining these masks with the previously complete phase domains
preserves all ten R928 and all fourteen R941 pattern allocations; each
has an exact rational feasible witness. For R928 the stronger oriented
$\Q(\zeta_{24})$ cone includes imaginary distinguished $STS$ rows,
both complex components of the central-$g$ $STS$ equations against
the distinguished objects, and every individual spectral fibre, with the
two endpoint-mixture parameters. All ten allocations still pass over
nonnegative rational variables.

An integer refinement reduces this R928 domain to four patterns. Each
orbit of the central semion $G$ has exactly one positive-charge
representative, since $G$ has charge $-1$. Associate to that orbit the
six normalized relative twist traces over $\Q(\zeta_{24})$. Relative
Galois averaging changes the trace atom but leaves its coefficient one;
therefore each of the $670$ phase-atom coefficients counts actual orbits
and must be a nonnegative integer. This applies also to the two
unrestricted restriction types, whose complete all-order trace domain
includes the zero tail. Rotation auxiliaries and both endpoint-mixture
parameters remain arbitrary nonnegative real variables.

After one-signed zero-equation deductions, complete binary branching on
$x\leq k$ or $x\geq k+1$ for an integer phase variable gives exact
infeasibility certificates for patterns $44,58,196,243,245,248$.
The respective trees have $13,13,11,13,31,11$ nodes and $49$ rational
Farkas leaves in total. Each leaf contradicts only necessary equations,
so the complete branch trees exclude these six patterns without relying
on a numerical solver's status. Patterns $40,49,186,191$ have explicit
nonnegative integer phase counts and exact rational auxiliary witnesses
satisfying every equation of this relaxation. The replay is
\path{r928_r941/integer_followup/verify.py} in the new evidence directory.
These four witnesses do not supply the missing modular rows, full center
fusion rules or associators; R928 itself remains unresolved.

For R941 an independent degree-two fusion-algebra check has span rank
fifteen and verifies $900$ physical/spectral Gram identities per
allocation. Of its $75$ generator-times-basis products, $55$ have
zero projection residual with forced nonnegative integral coordinates;
the remaining $20$ have nonzero residual, leaving closure open in all
fourteen allocations. These tests provide restrictions and witnesses,
not a new exclusion or a model for R928 or R941.

\subsection{The irrational dimension-ratio sector of R1768}
\label{subsec:creative-r1768}

Work in the pseudo-unitary Galois form supplied by the earlier general
sphericalization theorem. Put $q=\sqrt{37}$, $u=6+q$ and
$D=444+72q$. The source dimensions in canonical order are
\[
 (1,1,5+q,6+q,7+q,7+q,6+q,5+q).
\]
The saturated quadratic character has $h=-1/\sqrt{37}$, and
\[
 c_\chi(Z)=\frac{\sqrt{37}}{222}(37b_Z-6a_Z)
       \quad\text{for }\dim Z=a_Z+b_Z\sqrt{37}.
\]
Equation~\eqref{eq:creative-prime-masses} gives capacity $72$ at twist
one and capacity $144$ at each of the eighteen nonresidue primitive
37th twists, measured by the sum of $(37b_Z-6a_Z)^2$.
The unit and its Galois partner of dimension $\alpha=73+12q$ already
exhaust the twist-one capacity. Every other simple with irrational
$\dim Z/D$ therefore has one of those eighteen primitive twists.

\begin{proposition}[Complete irrational-sector restriction]
\label{prop:creative-r1768-sector}
Every pseudo-unitary Galois form of a categorification of R1768 has all
six noninvertible second indicators equal to $+1$.
Put $z=(0,0,1,1,1,1,1,1)$. Its center simples with irrational $\dim Z/D$, other than the unit
and its distinguished Galois partner, have exactly one of the following
forms:
\begin{enumerate}
\item $72$ simples of dimension $6u$, each with restriction $z$,
      four at every allowed twist;
\item $18$ simples of dimension $12u$, each with restriction $2z$,
      one at every allowed twist.
\end{enumerate}
The distinguished dimension-$\alpha$ object has the unique restriction
$(1,0,2,2,2,2,2,2)$.
\end{proposition}
\begin{proof}
In the previously complete high-column domain, the nonrational-ratio
dimensions avoiding the pointed object $g$ are $ku$ or $k\alpha$.
The latter cannot have primitive-37 twists: modular Galois symmetry
sends its dimension to $k$, and an integer-dimensional central object
forgets entirely to $1,g$. Such an object has twist $1$ or $\pm i$,
respectively, by induction and the already forced nontrivial pointed
associator. Galois square action preserves twist order.

For dimension $ku$, the Galois sign is negative. All six rational
source-character rows of the normalized $S$-matrix therefore negate
their values on the restriction vector, while both FP-character values
are unchanged. Character orthogonality gives the exact complement
\[
 v\longmapsto\widehat v=(0,0,k/3-v_2,\ldots,k/3-v_7).
\]
Applying integrality and nonnegativity of this complement to the
complete dimension domain leaves just
\[
 v=(0,0,0,0,1,0,1,1),\quad
 w=(0,0,1,1,0,1,0,0),\quad 2v,\quad2w,\quad z,\quad2z.
\]
The dimensions of $v,w$ are $3u$, those of $2v,2w,z$ are $6u$,
and that of $2z$ is $12u$.

Galois covariance across all eighteen twists forces aggregate
multiplicities $9a,9a,9b,9b,18c,18d$, respectively. The fibre
capacities imply $a+4b+4c+16d=16$. Thus their entire Gram contribution
is one of the five matrices
\[
 G_{37}=9U(vv^{\mathsf T}+ww^{\mathsf T})+18Vzz^{\mathsf T},
 \qquad U+4V=16,\quad V=0,1,2,3,4.
\]
This aggregate statement also covers Galois orbits larger than
eighteen. The common primitive trace ray fixes all retained indicator,
rotation and atomic modular rows as well. Subtract each of these five
contributions from the exact previously complete cone and retain its
full all-order trace domain for the rational-ratio columns.

There are eight remaining second-indicator assignments, hence forty
branches. The exact certificates reject thirty-nine. The only surviving
branch has $V=4,U=0$ and indicators
$(-1,1,1,1,1,1,1)$ on $(g,X_2,\ldots,X_7)$.
The nonnegative integer capacity equation then has just
$(a,b,c,d)=(0,0,4,0),(0,0,0,1)$, proving the two stated possibilities.

Finally the Galois image of the unit has positive sign. All rational
character values of its restriction are one, and its two FP values are
$\alpha,\alpha'$. Orthogonality gives uniquely
\[
 e_0+\frac{\alpha-1}{D}\mathbf d
       +\frac{\alpha'-1}{D'}\mathbf d'
   =e_0+\frac{\mathbf d-\mathbf d'}{\sqrt{37}}
   =(1,0,2,2,2,2,2,2).
\]
\end{proof}

Even after this unique restriction is imposed, the surviving cone has an
exact rational witness. This is not an integral center construction.
Its remaining rational sector has dimensions $tD/12$ with positive
integral $t$ and $\sum t^2=70$; removing its known $I(1)$ constituents
leaves $\sum t^2=52$ for the high rational sector.
The complete forty-branch certificate and both rational witnesses are
replayed by \path{r1768/verify.py} in the new evidence directory.
Rational-ratio simples may also have primitive 37th twists; the
preceding counts do not count those objects, and their all-order domain
is retained. R1768 remains unresolved.

\subsection{A forced four-object packet in the R1154 center}
\label{subsec:creative-r1154-packet}

For R1154 put $D=540+240\sqrt5$ and use the positive spherical
normalization. In all three remaining center patterns, the four simple
constituents of $I(g)$ have dimension $D/4$. Label them so that their
twists are $(i,i,-i,-i)$, as forced by
$\sum\theta_j=0$ and $\sum\theta_j^2=-4$.

\begin{proposition}[Semionic packet reconstruction]
\label{prop:creative-r1154-packet}
The constituents satisfy $Z_0^*=Z_1$, $Z_2^*=Z_3$, and their normalized
$S$ block is
\[
 B=\begin{pmatrix}
 -\frac14+\frac i2&-\frac14-\frac i2&\frac14&\frac14\\
 -\frac14-\frac i2&-\frac14+\frac i2&\frac14&\frac14\\
 \frac14&\frac14&-\frac14-\frac i2&-\frac14+\frac i2\\
 \frac14&\frac14&-\frac14+\frac i2&-\frac14-\frac i2
 \end{pmatrix}.
\]
Their four full normalized $S$ columns and all four fusion matrices are uniquely
determined by the previously retained restriction data. Every one of
the twenty-five pair histograms passes the resulting exact integrality
and nonnegativity tests.
\end{proposition}
\begin{proof}
Write $s=S/D$, let $Q$ project onto the four packet coordinates, and
put
\[
\begin{aligned}
 u&=\tfrac12(e_{Z_0}+e_{Z_1}+e_{Z_2}+e_{Z_3}),\\
 v&=\tfrac12(e_{Z_0}+e_{Z_1}-e_{Z_2}-e_{Z_3}),\\
 w&=\tfrac12(e_{A_0}+e_{A_1}-e_{A_2}-e_{A_3}).
\end{aligned}
\]
Here $A_0=1,A_1,\ldots,A_4$ are the distinguished twist-one objects.
For each retained histogram the spectral projector of the forced
$A_2$ matrix onto eigenvalue $-1$ is exactly
\[
 P=\frac{N_2(N_2-3)(N_2+5)(N_2^2-270N_2+225)}{7936}
   =ww^{\mathsf T}+Q-uu^{\mathsf T}.
\]
Verlinde diagonalization gives $P=sQs^*$, whence $su=w$ and
$s(\operatorname{im}Q\cap u^\perp)=\operatorname{im}Q\cap u^\perp$.
Since $s^2$ is duality, $sw=u$; applying $sTs=T^{-1}sT^{-1}$ to
$w$, using $Tu=iv$, gives $sv=-v$.

On the remaining two-dimensional space with orthonormal basis
$(e_{Z_0}-e_{Z_1})/\sqrt2,(e_{Z_2}-e_{Z_3})/\sqrt2$, the complete
symmetric-unitary solutions of the modular relation compatible with
charge conjugation are
\[
 A=\operatorname{diag}(i,-i),\qquad
 A=\frac i2\begin{pmatrix}-1&\pm\sqrt3\\\pm\sqrt3&1\end{pmatrix}.
\]
Each has $A^2=-I$, forcing the two dual pairs. In either mixed
solution the $(Z_2,Z_3)$ entry of $N_{Z_0}-N_{Z_1}$ is
$\pm\sqrt3$, contradicting integrality. The diagonal solution gives
the displayed block. Outside the packet every one of these four
normalized $s$ columns equals $w/2$, so the full columns are known.

Let $M$ be the restriction matrix and $V$ the matrix of these four
full normalized $S$ columns. The projection formula gives
$J=N_{I(g)}=M^{\mathsf T}N_gM$. Set
\[
\begin{aligned}
 D_+&=V\operatorname{diag}(4(B_{0j}-B_{1j}))V^*,\\
 D_-&=V\operatorname{diag}(4(B_{2j}-B_{3j}))V^*,\\
 H&=V\operatorname{diag}(-4,-4,4,4)V^*.
\end{aligned}
\]
Then
\[
 N_{Z_0},N_{Z_1}=\frac{J+H\pm2D_+}{4},\qquad
 N_{Z_2},N_{Z_3}=\frac{J-H\pm2D_-}{4}.
\]
The exact replay checks all entries, vacuum rows, duality transpose,
commutation with each other and $N_2$, and the sum identity in all
twenty-five histograms.
\end{proof}

This packet is not asserted to be a fusion subcategory. The remainder
of the center and the source associators are still unknown.
A separate necessary-data witness for pattern $6971$, pair $1$,
triple $0$ extends all seventy words of degree at most four and passes
$4900$ Gram and $560$ restriction checks. It is one verified witness,
not an exhaustive quartic enumeration or a categorification.

All twenty-five packet branches also have full tensor support: for each
packet object $Z$, the vacuum row of
$1+N_Z+N_Z^{\mathsf T}+(N_Z+N_Z^{\mathsf T})^2$ is strictly positive.
Thus $Z$ and $Z^*$ generate every center label in at most two steps.
Nevertheless the known forced fusion operators leave a common
$(n-9)$-dimensional scalar eigenspace, and their joint complex
commutant has dimension $9+(n-9)^2$. This identifies remaining spectral
freedom without constructing the missing fusion matrices.

A different construction test eliminates a tempting dimension match.
There is an exact identity
\[
 \FPdim(SU(3)_7\boxtimes\mathrm{Fib})=5(540+240\sqrt5).
\]
In the full seventy-two-object product, every simple dimension has
nonnegative $\sqrt5$ coefficient, and this coefficient vanishes
precisely for the three invertibles. Hence a rational-dimensional
connected separable algebra is supported on those invertibles.
Each invertible occurs with multiplicity at most one: the regular
right algebra-module $A$ is simple because
$\operatorname{End}_A(A)=\operatorname{Hom}(1,A)=\C$; for invertible
$g$, the module $gA$ is also simple, so free-module adjunction and
Schur's lemma give $\dim\operatorname{Hom}(g,A)\leq1$.
Therefore no connected separable algebra has dimension five, excluding
this direct connected-\'{e}tale module-category route to R1154.
The argument generalizes whenever a rational-linear field functional
vanishes on $1$, is nonnegative on all simple FP dimensions, and is
strictly positive on every noninvertible. It does not exclude longer
constructions or R1154 itself. The exact dimension calculation is
recorded in \path{r1154/quantum_subgroup_route.json}.

\subsection{The complete based-ring restriction on the R982 Morita class}
\label{subsec:creative-r982-morita}

\begin{proposition}[Morita partners preserve the R982 fusion ring]
\label{prop:creative-r982-morita}
Every fusion category in the Morita class of a complex categorification
of R982 has based fusion ring R982. This has no bound on the rank or
multiplicity of a possible Morita partner and does not assert equivalence
of its tensor categories.
\end{proposition}
\begin{proof}
Work first in a pseudo-unitary Galois form and conjugate any module
associators along with the category. The previous complete calculation
gives three necessary indecomposable based-module types. Each is
isomorphic over $\Q$ to the regular representation, although two are
not regular as based modules.

Here is the relevant general full-center argument. Put
$A_C=I_C(1)$ and transport $A_M=I_{C_M^*}(1)$ into $Z(C)$ using the
module-induced center equivalence. The module trace exists in the
pseudo-unitary case. The full-center and character identifications of
Shimizu, \emph{Further Results on the Structure of (Co)Ends in Finite
Tensor Categories}, Theorems~5.4, 5.11 and Corollary~5.13
(\href{https://arxiv.org/pdf/1801.02493}{arXiv:1801.02493}), identify
the module representation with $\operatorname{Hom}(A_M,A_C)$ over
$\operatorname{End}(A_C)$.
If a central simple has multiplicity $e_j$ in $A_C$, regularity forces
its multiplicity in $A_M$ also to be $e_j$. Thus $A_C$ is a summand
of $A_M$. Their common FP dimension $\FPdim(C)$ eliminates a
complement. Therefore $A_M\simeq A_C$ as objects. Since the source
Grothendieck algebra is commutative of dimension eight, the dual also
has exactly eight simples and is commutative. Equality of algebra
structures on the Lagrangian objects is not assumed.

Write $M_i,P_j$ for the left and right action matrices of the Morita
bimodule, in row-action convention. Internal-Hom associativity gives
\[
 \sum_i(M_i)_{ab}(M_i)_{cd}
       =\sum_j(P_j)_{ac}(P_j)_{bd}.
\]
Let $U,V$ have columns the flattened $M_i,P_j$. For every one of
the three left module types, the exact Gram identity is
$VV^{\mathsf T}=UU^{\mathsf T}$, with $U$ of rank eight.
Hence $V=UQ$ for an orthogonal eight-by-eight matrix $Q$.
The row lattice of $U$ is $\Z^8$ for the regular type and
$D_8=\{x\in\Z^8:\sum_i x_i\equiv0\pmod2\}$ for each other
type. Explicit row minors of determinant one or two certify these
equalities.

In the regular case integrality of $V$ makes $Q$ an integral
orthogonal matrix. In the $D_8$ cases, it makes
$\Z^8Q^{\mathsf T}$ an odd integral unimodular index-two
overlattice of $D_8$. The three possible nonzero discriminant cosets
are represented by $e_1$, $h=(1/2,\ldots,1/2)$ and $h-e_1$.
The latter two give even lattices, so only the $\Z^8$ overlattice
is possible. Thus $Q$ is a signed permutation in every case.
Nonnegativity of $V$ removes negative signs. After relabelling, the
right actions are exactly the left actions. Their linear independence
recovers every source multiplication coefficient from matrix products,
proving the based-ring identity. Galois conjugation leaves all these
integers unchanged.
\end{proof}

The matrix replay verifies all $4096$ bimodule Gram identities for
each type and the integral row-lattice certificates. The lattice argument itself does not construct module associators.
The exceptional realization in Theorem~\ref{thm:creative-e21-model}
separately supplies an algebra structure on $1+X_5$ and both
nonregular based-module types; uniqueness of their algebra or module
structures is not claimed.

\subsection{A Galois Gram bound and the two R3472 module branches}
\label{subsec:creative-r3472-gram}

\begin{lemma}[Integral-dimension Galois Gram bound]
\label{lem:creative-galois-gram}
Let $R$ be a fusion ring with Frobenius--Perron dimensions $d_i$, put
$K=\Q(d_i:i)$ and $D=\sum_i d_i^2$, and let
$h=\sum_{d_i\in\Z}d_i^2$. Then
\[
 h\leq \frac1{\operatorname{Tr}_{K/\Q}(D^{-1})}.
\]
If equality holds, every nonintegral $d_i$ satisfies
$\operatorname{Tr}_{K/\Q}(d_i/D)=0$.
Neither categorification nor closure of the integer-dimensional basis
elements is required.
\end{lemma}
\begin{proof}
Every conjugate of $d_i^2$ is an eigenvalue of the positive semidefinite
integer matrix $N_iN_{i^*}=N_iN_i^{\mathsf T}$, so $K$ is totally real.
For distinct embeddings $\sigma,\tau:K\longrightarrow\mathbb R$, the
dimension characters $\sigma(d_i)$ and $\tau(d_i)$ are distinct.
Evaluating the $\tau$-dimension character on the regular element
$\sum_i\sigma(d_i)X_i$, which satisfies
$X_j\sum_i\sigma(d_i)X_i=\sigma(d_j)\sum_i\sigma(d_i)X_i$, gives
\[
 \sum_i\sigma(d_i)\tau(d_i)=0.
\]
The Gram matrix of the conjugate dimension rows is therefore
$\operatorname{diag}(\sigma(D))$. Its integer-dimensional columns
contribute $h\mathbf1\mathbf1^{\mathsf T}$. The residual Gram matrix
is positive semidefinite. Evaluate its quadratic form at
$v_\sigma=1/\sigma(D)$ and put
$S=\operatorname{Tr}_{K/\Q}(D^{-1})>0$ to obtain $S-hS^2\geq0$.
If equality holds, the same quadratic form is a sum of squares, one
for each nonintegral dimension; each square is
$\operatorname{Tr}_{K/\Q}(d_i/D)^2$ and hence vanishes.
\end{proof}

The next assertions are conditional on the indicated based module lifting
to an indecomposable module category. No existence of either lift is
asserted. Put $r=\sqrt2$ and $D=700+490r$. Up to interchanging the two
module simples, the two source action matrices are
\[
 \begin{pmatrix}7&7\\7&5\end{pmatrix},\qquad
 \begin{pmatrix}11&5\\5&1\end{pmatrix}.
\]
Their full integral commutants and positive dimension eigenvalues are
\[
 \begin{array}{c|c}
 \text{matrix}&\text{dimension}\\\hline
 \begin{pmatrix}s+t&7t\\7t&s-t\end{pmatrix}&s+5tr\\[3pt]
 \begin{pmatrix}s+t&t\\t&s-t\end{pmatrix}&s+tr
 \end{array}
 \qquad t\geq0,\quad s\geq t,\quad s,t\in\Z.
\]
Morita invariance gives global dimension $D$ for either dual category
\cite[Corollary~7.16.7]{EGNO}. The regular dimension element acts with
rank one on the module, as in \cite[Exercise~7.16.8]{EGNO}.

Lemma~\ref{lem:creative-galois-gram} gives $h\leq7$. Here the integral
simples form a fusion subcategory, since all radical coefficients are
nonnegative. By categorical Lagrange integrality
\cite[Theorem~7.17.6]{EGNO}, $D/h$ is integral, hence
$1400/h,9800/h^2\in\Z$ and $h\in\{1,2,5,7\}$.
The subcategory is pointed: the only nonpointed integral profiles to
consider are $1,2$ for $h=5$ and $1,1,1,2$ for $h=7$; the square of the
unique dimension-two simple would give $4=1+2a$ or $4=3+2a$,
respectively. Thus write $g=h$ for its cyclic group order.

\begin{proposition}[Three further first-branch exclusions]
\label{prop:creative-r3472-three}
For the first branch, the ten previously retained dual dimension spectra
reduce to seven. The following three spectra cannot be fusion rings:
\[
\begin{array}{c|l}
g&\text{nonintegral dimensions with multiplicity}\\\hline
1&(6+5r)^2,\ (9+5r)^1,\ (14+10r)^1\\
1&(5+5r)^1,\ (8+5r)^2,\ (14+10r)^1\\
2&(7+5r)^2,\ (9+5r)^1,\ (13+10r)^1.
\end{array}
\]
Superscripts in this table denote multiplicities, not powers.
\end{proposition}
\begin{proof}
In each row the repeated pair is either self-dual or exchanged by
duality; all other nonintegral dimensions occur once. For each of these
six possibilities introduce nonnegative integers
$m_{ijk}=[\one:X_iX_jX_k]$, identifying cyclic rotations and the
reverse-dual triple $(k^*,j^*,i^*)$. Unit triples and cyclic-group
products are fixed. Equating rational and radical coefficients in
every product dimension yields a finite nonnegative integer linear
system. No commutativity or associativity assumption is used.

The exact verifier exhausts all six systems using interval bounds,
row gcd congruences and complete finite branching. Their search-tree
sizes are $72,22,39,11,380,79$, and every tree is empty. Every variable
has a finite upper bound from a positive coefficient in a dimension
equation; there is no node or time cutoff. The original R3472 tensor
is checked as a positive control for the equation builder. Thus the
six systems supply exact incompatibility certificates for the three
spectra. The receipt retains the seven other spectra only as necessary
possibilities, without claiming fusion multiplications.
\end{proof}

For the second branch the regular-element identities give
\[
 \sum s_i^2=350,\qquad\sum t_i^2=175,\qquad\sum s_it_i=245.
\]
After removing the $g$ invertibles and putting $u_i=s_i-t_i$, the
complete arithmetic domain is
\[
 \sum u_i^2=35-g,\qquad\sum u_it_i=70,\qquad\sum t_i^2=175,
 \qquad u_i\geq0,\ t_i\geq1.
\]
In particular $u_i\leq5$ and $t_i\leq13$. A finite multiplicity
recursion over these types, completed by square partitions for $u=0$,
counts each dimension multiset once. The exact counts are
\[
 \begin{array}{c|rrrr|r}
 g&1&2&5&7&\text{total}\\\hline
 \text{spectra}&272348&124228&3900&2&400478.
 \end{array}
\]
When $g=7$, equality in Lemma~\ref{lem:creative-galois-gram} gives
$5s_i=7t_i$. Writing $\delta=7+5\sqrt2$, the two possibilities are
$1^7,\delta^7$ and $1^7,\delta^3,(2\delta)^1$.
Since $\delta^2=1+14\delta$, the stabilizer of a $\delta$-dimensional
simple is trivial and its orbit under $C_7$ has size seven. The second
possibility is impossible. Thus $400477$ arithmetic spectra survive;
full multiplication tests for this entire domain are not supplied.

The exact programs, source identities and receipts are in
\path{evidence/rank8_creative_20261022/r3472/}. Normal and optimized
Python replay produce identical results. These refinements are
conditional Morita restrictions. R3472 was unresolved at this checkpoint; its later general exclusion is Theorem~\ref{thm:r3472-general-exclusion}.

\subsection{A common cyclic-parent obstruction for R1768 and R3472}
\label{subsec:creative-cyclic-parent}

\begin{lemma}[Cyclic stabilizer divisibility]
\label{lem:creative-cyclic-stabilizer}
Let $H$ be a cyclic subgroup of invertibles in a complex fusion category
$\mathcal C$, and let $X$ be a simple object with $Xh\simeq X$ for all
$h\in H$. Suppose $H$ lifts to a Tannakian pointed subcategory of
$Z(\mathcal C)$ whose forgetful image in $\mathcal C$ is fully faithful.
Then $\FPdim(X)/|H|$ is an algebraic integer.
\end{lemma}
\begin{proof}
Let $A=\bigoplus_{h\in H}h$ be the regular algebra and
$\mathcal D=\mathcal C_A$ the ordinary de-equivariantization
\cite[Theorem~8.23.3]{EGNO}. Module structures supported on copies
of $X$ are the representations of a twisted group algebra of $H$.
Every scalar two-cocycle on a cyclic group over $\C$ is a coboundary:
rescale a projective generator with $T^{|H|}=\lambda$ by a root of
$\lambda$. Hence there are $|H|$ simple $A$-modules with underlying
object $X$. Their dimensions in $\mathcal D$ are
$\FPdim_{\mathcal C}(X)/\FPdim_{\mathcal C}(A)$.
For completeness, this normalization follows because the positive
dimension vectors before and after forgetting satisfy the same
indecomposable $\mathcal C$-module eigenvector equations; evaluating
at the unit $A$ of $\mathcal D$ fixes the constant to $|H|$.
Dimensions of fusion-category simples are algebraic integers.
\end{proof}

\begin{corollary}[Cyclic quotient square divisibility]
\label{cor:creative-ng-square}
An ordinary de-equivariantization of a category with near-group rules
$\mathrm{NG}(C_n,n)$ by a group of order $h$ necessarily has $h^2\mid n$.
\end{corollary}
\begin{proof}
Every simple in the fully faithful image of the de-equivariantizing
representation category has integral dimension, hence is invertible.
Its group of simple objects is therefore cyclic. The unique
noninvertible of the near-group category is fixed by this subgroup.
Its dimension $d$ satisfies $d^2=n+nd$ and has quadratic norm $-n$.
Lemma~\ref{lem:creative-cyclic-stabilizer} makes $d/h$ integral, so its
norm $-n/h^2$ belongs to $\Z$. The discriminant $n^2+4n$ lies strictly
between $(n+1)^2$ and $(n+2)^2$, so this is a genuine quadratic norm.
\end{proof}

Set $\delta=g+\sqrt{g^2+1}$ and $D_g=g(1+\delta^2)$.
The hypothetical cyclic near-group parent $\mathrm{NG}(C_{4g^2},4g^2)$
has noninvertible dimension $2g\delta$ and global dimension $4gD_g$.
Dividing its global dimension by $4g$ would give the desired
multiplicity-two generalized-near-group scalar profile, but
$(4g)^2\nmid4g^2$. Equivalently, the quotient would contain a simple
of dimension $\delta/2$, whose norm is $-1/4$.
For $g=6$ and $g=7$ this excludes, respectively, the direct ordinary
quotient routes
\[
 \mathrm{NG}(C_{144},144)\longrightarrow D=444+72\sqrt{37},
 \qquad
 \mathrm{NG}(C_{196},196)\longrightarrow D=700+490\sqrt2.
\]
No assertion that these parents exist is required. The obstruction
does not exclude noncyclic parents, other Morita representatives or
generalized-Haagerup parents. In particular the scalar possibilities
of a $C_{12}$ or $C_{14}$ generalized-Haagerup parent followed by a
$C_2$ quotient are not settled here. These are restrictions on
construction routes for R1768 and R3472, not general exclusions of
those rings or constructions of their associators.

